\begin{table*}[!ht]
\centering
\caption{Held-out performance of the prespecified L2-regularized logistic-regression ablations. AUC confidence intervals use 1,000 class-stratified bootstrap samples.}
\label{tab:model_performance}
\small
\setlength{\tabcolsep}{5pt}
\resizebox{\textwidth}{!}{%
\begin{tabular}{lrrrrrr}
\toprule
Feature configuration & Test n & AUC (95\% CI) & Balanced accuracy & Sensitivity & Specificity & Brier score \\
\midrule
Clinical core & 141 & 0.580 (0.486--0.674) & 0.596 & 0.804 & 0.389 & 0.230 \\
Clinical extended & 141 & 0.608 (0.515--0.699) & 0.537 & 0.941 & 0.133 & 0.225 \\
Clinical core + NfL & 141 & 0.640 (0.553--0.731) & 0.649 & 0.843 & 0.456 & 0.213 \\
Clinical extended + NfL & 141 & 0.678 (0.590--0.762) & 0.636 & 0.706 & 0.567 & 0.210 \\
Clinical extended + DTI & 141 & 0.636 (0.547--0.727) & 0.593 & 0.765 & 0.422 & 0.221 \\
Clinical extended + standard-mean DTI & 141 & 0.628 (0.533--0.721) & 0.596 & 0.804 & 0.389 & 0.222 \\
Clinical extended + FreeSurfer* & 141 & 0.780 (0.703--0.848) & 0.678 & 0.667 & 0.689 & 0.178 \\
All modalities* & 141 & 0.802 (0.730--0.872) & 0.689 & 0.588 & 0.789 & 0.176 \\
\bottomrule
\end{tabular}
}
\begin{minipage}{0.98\linewidth}\footnotesize
*Exploratory structural-MRI analyses were not part of the original clinical/NfL/DTI analysis scope. Classification thresholds were selected from training-only out-of-fold predictions.
\end{minipage}
\end{table*}